%% Separately attributed additions retained after exact source restoration.
\section*{Editorial supplement to the restored chapter}
\addcontentsline{toc}{section}{Editorial supplement to the restored chapter}
\noindent\textit{Reader note.} The following two results are separately
authored additions retained from the historical draft.  The
Keisler--Shelah statement is controlled by its original source, and the
compactness theorem is accompanied by a complete ultraproduct proof in the
Part~III annex.  Neither result alters or replaces the exact inherited chapter
above \cite{Shelah1971Ultrapowers,ChangKeisler1990}.

\begin{hbtheorem}{thm:p03-keisler-shelah}{cited=Shelah1971Ultrapowers}{Keisler--Shelah isomorphism theorem}
For structures \(\mathcal M\) and \(\mathcal N\) in the same signature,
\(\mathcal M\equiv\mathcal N\) if and only if some ultrapower of
\(\mathcal M\) is isomorphic to some ultrapower of \(\mathcal N\).
No fixed index set, ultrafilter, nonprincipality condition, or cardinal bound
is asserted.  The reverse implication follows from \L o\'s's theorem and
isomorphism invariance; the forward implication is Shelah's theorem.
\end{hbtheorem}

\begin{hbtheorem}{thm:p03-compactness-ultraproduct}{proved=annex:p3.types.8}{Compactness via ultraproducts}
Assume the Axiom of Choice in the ambient set theory.  If every finite subset
of a set \(\Gamma\) of first-order \(\Lang\)-sentences is satisfiable, then
\(\Gamma\) is satisfiable.  The derivation in the annex uses the finite-subset
index set, extends its generated proper filter by the Ultrafilter Lemma, and
then applies \L o\'s's theorem; no choice-free claim is made.
\end{hbtheorem}
